%%%%%%%%%%%%%%%%%%%%%%%%%%%%%%%%%%%%%%%%%
% Cleese Assignment (For Students)
% LaTeX Template
% Version 2.0 (27/5/2018)
%
% This template originates from:
% http://www.LaTeXTemplates.com
%
% Author:
% Vel (vel@LaTeXTemplates.com)
%
% License:
% CC BY-NC-SA 3.0 (http://creativecommons.org/licenses/by-nc-sa/3.0/)
% 
%%%%%%%%%%%%%%%%%%%%%%%%%%%%%%%%%%%%%%%%%

%----------------------------------------------------------------------------------------
%	PACKAGES AND OTHER DOCUMENT CONFIGURATIONS
%----------------------------------------------------------------------------------------

\documentclass[11pt]{article}

%%%%%%%%%%%%%%%%%%%%%%%%%%%%%%%%%%%%%%%%%
% Cleese Assignment
% Structure Specification File
% Version 1.0 (27/5/2018)
%
% This template originates from:
% http://www.LaTeXTemplates.com
%
% Author:
% Vel (vel@LaTeXTemplates.com)
%
% License:
% CC BY-NC-SA 3.0 (http://creativecommons.org/licenses/by-nc-sa/3.0/)
% 
%%%%%%%%%%%%%%%%%%%%%%%%%%%%%%%%%%%%%%%%%

%----------------------------------------------------------------------------------------
%	PACKAGES AND OTHER DOCUMENT CONFIGURATIONS
%----------------------------------------------------------------------------------------

\usepackage{lastpage} % Required to determine the last page number for the footer

\usepackage{graphicx} % Required to insert images

\setlength\parindent{0pt} % Removes all indentation from paragraphs

\usepackage[most]{tcolorbox} % Required for boxes that split across pages

\usepackage{booktabs} % Required for better horizontal rules in tables

\usepackage{listings} % Required for insertion of code

\usepackage{etoolbox} % Required for if statements

\usepackage[T1]{fontenc} % Output font encoding for international characters

\usepackage[utf8]{inputenc} % Required for inputting international characters

%----------------------------------------------------------------------------------------
%	MARGINS
%----------------------------------------------------------------------------------------

\usepackage{geometry} % Required for adjusting page dimensions and margins

\geometry{
	paper=a4paper, % Change to letterpaper for US letter
	top=3cm, % Top margin
	bottom=3cm, % Bottom margin
	left=2.5cm, % Left margin
	right=2.5cm, % Right margin
	headheight=14pt, % Header height
	footskip=1.4cm, % Space from the bottom margin to the baseline of the footer
	headsep=1.2cm, % Space from the top margin to the baseline of the header
	%showframe, % Uncomment to show how the type block is set on the page
}

%----------------------------------------------------------------------------------------
%	FONT
%----------------------------------------------------------------------------------------

\usepackage[utf8]{inputenc} % Required for inputting international characters
\usepackage[T1]{fontenc} % Output font encoding for international characters

\usepackage[sfdefault,light]{roboto} % Use the Roboto font

%----------------------------------------------------------------------------------------
%	HEADERS AND FOOTERS
%----------------------------------------------------------------------------------------

\usepackage{fancyhdr} % Required for customising headers and footers

\pagestyle{fancy} % Enable custom headers and footers

\lhead{\small\assignmentClass{ : }{}\ \assignmentTitle} % Left header; output the instructor in brackets if one was set
\chead{} % Centre header
\rhead{\small\ifdef{\assignmentAuthorName}{\assignmentAuthorName}{\ifdef{\assignmentDueDate}{Due\ \assignmentDueDate}{}}} % Right header; output the author name if one was set, otherwise the due date if that was set

\lfoot{} % Left footer
\cfoot{\small Page\ \thepage\ of\ \pageref{LastPage}} % Centre footer
\rfoot{} % Right footer

\renewcommand\headrulewidth{0.5pt} % Thickness of the header rule

%----------------------------------------------------------------------------------------
%	MODIFY SECTION STYLES
%----------------------------------------------------------------------------------------

\usepackage{titlesec} % Required for modifying sections

%------------------------------------------------
% Section

\titleformat
{\section} % Section type being modified
[block] % Shape type, can be: hang, block, display, runin, leftmargin, rightmargin, drop, wrap, frame
{\Large\bfseries} % Format of the whole section
{\thesection} % Format of the section label
{6pt} % Space between the title and label
{} % Code before the label

\titlespacing{\section}{0pt}{0.5\baselineskip}{0.5\baselineskip} % Spacing around section titles, the order is: left, before and after

%------------------------------------------------
% Subsection

\titleformat
{\subsection} % Section type being modified
[block] % Shape type, can be: hang, block, display, runin, leftmargin, rightmargin, drop, wrap, frame
{\itshape} % Format of the whole section
{(\alph{subsection})} % Format of the section label
{4pt} % Space between the title and label
{} % Code before the label

\titlespacing{\subsection}{0pt}{0.5\baselineskip}{0.5\baselineskip} % Spacing around section titles, the order is: left, before and after

\renewcommand\thesubsection{(\alph{subsection})}

%----------------------------------------------------------------------------------------
%	CUSTOM QUESTION COMMANDS/ENVIRONMENTS
%----------------------------------------------------------------------------------------

% Environment to be used for each question in the assignment
\newenvironment{question}{
	\vspace{0.5\baselineskip} % Whitespace before the question
	\section{} % Blank section title (e.g. just Question 2)
	\lfoot{\small\itshape\assignmentQuestionName~\thesection~continued on next page\ldots} % Set the left footer to state the question continues on the next page, this is reset to nothing if it doesn't (below)
}{
	\lfoot{} % Reset the left footer to nothing if the current question does not continue on the next page
}

%------------------------------------------------

% Environment for subquestions, takes 1 argument - the name of the section
\newenvironment{subquestion}[1]{
	\subsection{#1}
}{
}

%------------------------------------------------

% Command to print a question sentence
\newcommand{\questiontext}[1]{
	\textbf{#1}
	\vspace{0.5\baselineskip} % Whitespace afterwards
}

%------------------------------------------------

% Command to print a box that breaks across pages with the question answer
\newcommand{\answer}[1]{
	\begin{tcolorbox}[breakable, enhanced]
		#1
	\end{tcolorbox}
}

%------------------------------------------------

% Command to print a box that breaks across pages with the space for a student to answer
\newcommand{\answerbox}[1]{
	\begin{tcolorbox}[breakable, enhanced]
		\vphantom{L}\vspace{\numexpr #1-1\relax\baselineskip} % \vphantom{L} to provide a typesetting strut with a height for the line, \numexpr to subtract user input by 1 to make it 0-based as this command is
	\end{tcolorbox}
}

%------------------------------------------------

% Command to print an assignment section title to split an assignment into major parts
\newcommand{\assignmentSection}[1]{
	{
		\centering % Centre the section title
		\vspace{2\baselineskip} % Whitespace before the entire section title
		
		\rule{0.9\textwidth}{0.5pt} % Horizontal rule
		
		\vspace{0.75\baselineskip} % Whitespace before the section title
		{\LARGE \MakeUppercase{#1}} % Section title, forced to be uppercase
		
		\rule{0.9\textwidth}{0.5pt} % Horizontal rule
		
		\vspace{\baselineskip} % Whitespace after the entire section title
		\setcounter{section}{0}
	}
}

%----------------------------------------------------------------------------------------
%	TITLE PAGE
%----------------------------------------------------------------------------------------


\author{\textbf{\assignmentAuthorName}} % Set the default title page author field
\date{} % Don't use the default title page date field

\title{
	\textsc{\LARGE Haute École d'Ingénierie et de Gestion du Canton de Vaud}\\[1.5cm] % logo Ecole
	\includegraphics[scale=.2]{images/heig.png}\\[1cm] % Include a department/university logo - this will require the graphicx package
	\thispagestyle{empty} % Suppress headers and footers
	\vspace{0.1\textheight} % Whitespace before the title
	\textbf{\assignmentTitle}\\[4pt]
	\ifdef{\assignmentDueDate}{{\assignmentDueDate}\\}{} % If a due date is supplied, output it
	\ifdef{\assignmentClassInstructor}{{\large \textit{\assignmentClassInstructor}}}{} % If an instructor is supplied, output it
	\vspace{0.32\textheight} % Whitespace before the author name
}
\setcounter{page}{0} % Include the file specifying the document structure and custom commands

%----------------------------------------------------------------------------------------
%	ASSIGNMENT INFORMATION
%----------------------------------------------------------------------------------------

% Required
\newcommand{\assignmentQuestionName}{Question} % The word to be used as a prefix to question numbers; example alternatives: Problem, Exercise
\newcommand{\assignmentClass}{CNM} % Course/class
\newcommand{\assignmentTitle}{Acceleration of an AI application} % Assignment title or name
\newcommand{\assignmentAuthorName}{Kevin Jordil \& Olivier D'Ancona} % Student name


% Optional (comment lines to remove)
\newcommand{\assignmentClassInstructor}{Professor: Marina Zapater } % Intructor name/time/description
\newcommand{\assignmentDueDate}{Thursday,\ 02\ February\, 2023} % Due date

%----------------------------------------------------------------------------------------

\begin{document}

%----------------------------------------------------------------------------------------
%	TITLE PAGE
%----------------------------------------------------------------------------------------

\maketitle % Print the title page

\thispagestyle{empty} % Suppress headers and footers on the title page

\newpage

\assignmentSection{Stage 1 – Choosing an application}

\section{Multi-Layer Perceptron}

\paragraph*{Multi-Layer Perceptron (MLP)}is a type of artificial neural network commonly used for supervised learning tasks such as classification.
It consists of one input layer, some hidden layer and one output layer. In each layer, each neuron is connected to all neurons in the next layer.
The architecture that we built have 3 layers:

\paragraph*{Architecture} our neural network implementation is as follows:

\begin{itemize}
	\item Input layer $I$ of size 784 (28x28 images)
	\item Hidden layer $H$ of size 1000 (number of hidden neurons)
	\item Output layer $O$ of size 10 (number of class of the dataset)
\end{itemize}

We used a c implementation who consists of a single file with all the functions needed to train and test the neural network. The batch size $B=10$.

\paragraph*{Dataset} We used the fashion dataset which is a commonly used dataset in machine learning and computer vision task.
He consists of a collection of Zalando shopping items such as shirts, pants and shoes. There are in total 10 categories.
he class are T-shirt/top, Trouser, Pullover, Dress, Coat, Sandal, Shirt, Sneaker, Bag and Ankle boot.
Each image is 28 pixels in height and 28 pixels in width, for a total of 784 pixels in total.
Each pixel has a single pixel-value associated with it, indicating the lightness or darkness of that pixel, with higher numbers meaning darker.
This pixel-value is an integer between 0 and 255. The training and test data sets have 785 columns.
The first column consists of the class labels (see above), and represents the article of clothing.
The rest of the columns contain the pixel-values of the associated image.

\section{Grayscale Conversion Program}

\paragraph*{Definition} We ran into troubles with the cuda neural network acceleration so we developped an auxiliary program to be accelerated with cuda.
The program converts a color image to a grayscale image which is a good idea to implement with our multi layer perceptron.
Because our mlp necessitates a grayscale image as input we can convert rgb images to grayscale images with this program.
Furthermore this task is a good example to show the acceleration of a program with cuda.

\paragraph*{Principle} For each pixel of the image, we calculate the average of all color channels.
The result is the grayscale value of the pixel. Then we replace the color channels with the grayscale value in a new image.

\pagebreak

\assignmentSection{Stage 2 – Analysing application bottlenecks}

\section{Execution time}

\paragraph*{MLP baseline} The baseline execution time is showed on table \ref{tab:mlp_baseline} with a total time of 575.676547s.

\begin{table}[h]	\centering
	\begin{tabular}{|c|c|c|c|c|c|}
		\hline
		\textbf{time} & \textbf{accuracy} & \textbf{error} & \textbf{samples/sec} & \textbf{gflop/s} \\ \hline
		67.270 s      & 40.99\%           & 0.045          & 743.28               & 2.20             \\ \hline
		134.516 s     & 61.63\%           & 0.022          & 743.53               & 2.20             \\ \hline
		201.754 s     & 74.50\%           & 0.019          & 743.64               & 2.20             \\ \hline
		268.999 s     & 82.60\%           & 0.016          & 743.54               & 2.20             \\ \hline
		336.225 s     & 87.68\%           & 0.015          & 743.77               & 2.20             \\ \hline
		403.459 s     & 90.91\%           & 0.013          & 743.67               & 2.20             \\ \hline
		470.688 s     & 93.01\%           & 0.013          & 743.73               & 2.20             \\ \hline
		537.911 s     & 94.43\%           & 0.012          & 743.79               & 2.20             \\ \hline
	\end{tabular}
	\caption{MLP Baseline execution time}
	\label{tab:mlp_baseline}
\end{table}

\paragraph*{MLP accelerated with OpenMP} The accelerated execution time is showed on table \ref{tab:mlp_accelerated} with a total time of 393.993599s.

\begin{table}[h]	\centering
	\begin{tabular}{|c|c|c|c|c|c|}
		\hline
		\textbf{time} & \textbf{accuracy} & \textbf{error} & \textbf{samples/sec} & \textbf{gflop/s} \\ \hline
		46.092 s      & 40.99\%           & 0.045          & 1084.79              & 3.21             \\ \hline
		92.241 s      & 61.63\%           & 0.022          & 1083.43              & 3.20             \\ \hline
		138.231 s     & 74.50\%           & 0.019          & 1087.20              & 3.22             \\ \hline
		184.198 s     & 82.59\%           & 0.016          & 1087.76              & 3.22             \\ \hline
		230.153 s     & 87.68\%           & 0.015          & 1088.01              & 3.22             \\ \hline
		276.135 s     & 90.91\%           & 0.013          & 1087.40              & 3.22             \\ \hline
		322.124 s     & 93.01\%           & 0.013          & 1087.23              & 3.22             \\ \hline
		368.109 s     & 94.43\%           & 0.012          & 1087.32              & 3.22             \\ \hline
	\end{tabular}
	\caption{MLP OpenMP accelerated execution time}
	\label{tab:mlp_accelerated}
\end{table}

\paragraph*{MLP accelerated with OpenMP and CUDA} The accelerated execution time is showed on table \ref{tab:mlp_accelerated_cuda} with a total time of 326.487s.

\begin{table}[h]	\centering
	\begin{tabular}{|c|c|c|c|c|c|}
		\hline
		\textbf{time} & \textbf{accuracy} & \textbf{error} & \textbf{samples/sec} & \textbf{gflop/s} \\ \hline
		64.577 s      & 36.72\%           & 0.061          & 774.27               & 2.29             \\ \hline
		130.425 s     & 57.61\%           & 0.034          & 759.32               & 2.25             \\ \hline
		195.328 s     & 70.85\%           & 0.029          & 770.39               & 2.28             \\ \hline
		261.072 s     & 79.38\%           & 0.025          & 760.52               & 2.25             \\ \hline
		326.487 s     & 84.74\%           & 0.024          & 764.35               & 2.26             \\ \hline
		391.392 s     & 88.17\%           & 0.021          & 770.37               & 2.28             \\ \hline
		455.994 s     & 90.38\%           & 0.021          & 773.96               & 2.29             \\ \hline
		519.951 s     & 91.87\%           & 0.019          & 781.79               & 2.31             \\ \hline
		584.651 s     & 92.86\%           & 0.020          & 772.79               & 2.29             \\ \hline
		649.924 s     & 93.56\%           & 0.018          & 766.02               & 2.27             \\ \hline
		714.872 s     & 94.03\%           & 0.018          & 769.84               & 2.28             \\ \hline
		780.575 s     & 94.48\%           & 0.017          & 761.00               & 2.25             \\ \hline
		845.612 s     & 94.74\%           & 0.016          & 768.80               & 2.27             \\ \hline
	\end{tabular}
	\caption{MLP OpenMP and CUDA accelerated execution time}
	\label{tab:mlp_accelerated_cuda}
\end{table}

\paragraph*{Grayscale baseline} In the following results, the image name contains the resolution in pixels. Here is the baseline execution time:

\begin{itemize}
	\item ./main chicky\_512\_512.png -> Time elapsed: 3.431923 ms
	\item ./main avatar\_5000\_2381.jpg -> Time elapsed: 151.122006 ms
	\item ./main winter\_10667\_6000.jpg -> Time elapsed: 803.954043 ms
\end{itemize}

\paragraph*{Grayscale OpenMP} The following results use the same images but with CUDA acceleration:

\begin{itemize}
	\item	./main chicky\_512\_512.png -> Time elapsed: 2.095706 ms
	\item	./main avatar\_5000\_2381.jpg -> Time elapsed: 61.278724 ms
	\item	./main winter\_10667\_6000.jpg -> Time elapsed: 209.700440 ms
\end{itemize}

\section{Complexity Analysis}

\paragraph*{MLP} The MLP is a feedforward neural network, so the complexity is bounded by the biggest loop, which is the backpropagation loop in $O(\#epoch  \cdot H \cdot O \cdot B)$

\paragraph*{Grayscale} The complexity is one computation for every pixel so $O(m \cdot n )$ where $m$ is the number of rows and $n$ is the number of columns of the image.

\section{Theoretical acceleration}

\paragraph*{Hardware} The theoretical limitations depends on the hardware.
In our case, we use a jetson nano which has a max threads/block bound of 1024 and 128 CUDA cores.
Therefore, if we want to accelerate our application properly, we need want to maximize the number of thread running the same operation on the warps.
In our case, we have a lot of small operations in parallel and we therefore we can't use the full power of the GPU because memory transfer will be too long.

If we analyze the GPU time only with nprof on table \ref{tab:mlp_cuda_profiler}, we can use Amdahl's law to calculate the theoretical acceleration.
The kernel use 11.69 second, and the memcpy use 15.41 second and the memset use 5.6s. The part we can accelerate is only over those 11.69seconds.
Therefore, because the total GPU time takes more than running on the CPU baseline, we can't accelerate the application with the GPU faster than the cpu.

\paragraph*{Roofline} 

The max performance is 472 GFLOPs and the max memory bandwidth is 25.6G GB/s. Therefore, the jetson is balanced with CI = 18.43.
If we calculate the number of operations on the backpropagation kernel we have: $Y * H * B * 18op = 1.8M$. The number of data transfered is $Y * H * B * 4Bytes = 0.4M$. 
Therefore the CI of our kernel is 4.5. Using the symetry in the roofline model our kernel performance is $ 472 / 18.43 * 4.5 = 115.24 GFLOPs$.
As we can see, we are not limited by the performance but by the memory.

\assignmentSection{Stage 3 – Acceleration}

\section{Process}
First, we added OpenMP to the base code. As a result, we have a slight performance increase.
Then we wanted to add a CUDA kernel function as requested.
However, the performance was not as good as the original code.
So we spent a lot of time trying to optimize the CUDA code.
We tried to make several CUDA kernel functions, to use cuBLAS, to make memory optimizations.
Unfortunately, we never managed to improve the code base with CUDA
because the transfer between the CPU and the GPU took more time than the CUDA kernel function itself.
Therefore, we decided to make a second application and optimize it with a CUDA kernel function.

\section{Accelerated parts}

\paragraph*{Choice}
We chose to keep the optimization only with OpenMP on MLP because it is more efficient alone than with CUDA optimizations.
To respect the instructions and accelerate part of the code on GPU, we used the CUDA runtime library with CUDA kernel functions which allowed us to have good results.
We looked at using cuBLAS but it was not adapted because it is intended for matrix calculations and it is not really what is done in our example.
We did not use cuDNN because it is not suitable for our needs as it is specifically created for deep neural networks.

\paragraph*{Data} For MLP, we use the Fashion MINST dataset as explained above.
So we will have several large arrays but mostly loops that will go through these arrays.
So we found three triple for loops that took a lot of execution time according to us.
In a first step, we regenerated a single triple loop with these three loops, which will simplify the transition to CUDA kernel. In a second step, we put all the data on the GPU.
Then, we created a CUDA kernel function which allows to replace the generated triple for loop.


\begin{lstlisting}[language=C, caption="Original code with triple loop"]
	/* dy */
	for (int b = 0; b < B; b++)
	  for (int k = 0; k < Y; k++)
		dy[b * Y + k] = p[b * Y + k] - t[b * Y + k];
	/* dv := h * dy' */
	for (int b = 0; b < B; b++)
	  for (int j = 0; j < H; j++)
		for (int k = 0; k < Y; k++)
		  dv[k * H + j] += h[b * H + j] * dy[b * Y + k];
	/* dh := v * dy */
	for (int b = 0; b < B; b++)
	  for (int j = 0; j < H; j++)
		for (int k = 0; k < Y; k++)
		  dh[b * H + j] += v[k * H + j] * dy[b * Y + k];
\end{lstlisting}


\begin{lstlisting}[language=C, caption="Edited nested triple loop"]
	/* dy */
	for (int b = 0; b < B; b++)
	  for (int k = 0; k < Y; k++)
		for (int j = 0; j < H; j++)
		{
		  dy[b * Y + k] = p[b * Y + k] - t[b * Y + k];
		  dv[k * H + j] += h[b * H + j] * dy[b * Y + k];
		  dh[b * H + j] += v[k * H + j] * dy[b * Y + k];
		}
\end{lstlisting}

\begin{lstlisting}[language=C, caption="CUDA kernel function to calculate back\_propagation"]
	// Cuda kernel function of loop
	__global__ void backprop_kernel(float *p, float *t, float *dv, float *v, float *dh, float *h)
	{
		int y_idx = threadIdx.x + blockIdx.x * blockDim.x;
		int h_idx = threadIdx.y + blockIdx.y * blockDim.y;
	
		if (y_idx >= H || h_idx >= Y)
			return;
	
		float dy, dh_gpu, dv_gpu;
	
		for(int b_idx = 0; b_idx < B; b_idx++){
		  dy = p[b_idx * Y + h_idx] - t[b_idx * Y + h_idx];
		  dv_gpu = h[b_idx * H + y_idx] * dy;
		  dh_gpu = v[h_idx * H + y_idx] * dy;
		  atomicAdd(&dv[h_idx * H + y_idx], dv_gpu);
		  atomicAdd(&dh[b_idx * H + y_idx], dh_gpu);
		}
	}
	
	dim3 dimBlock(32, 32, 1);
	dim3 dimGrid(Y / 32 + 1, H / 32 + 1, 1);
	backprop_kernel<<<dimGrid, dimBlock>>>(p_gpu, t_gpu, dv_gpu, v_gpu, dh_gpu, h_gpu);
\end{lstlisting}

As said before, this code is still slower than the original and using OpenMP alone is faster.
So there is an OpenMP directive (\#pragma omp parallel for) before each for loop.
The number of threads is set to 4.

For the application that transforms the image into grayscale, there is a double for loop that runs through the image.
So we created a CUDA kernel function that allows to replace the generated double for loop and thus to use the GPU computing power.

\begin{lstlisting}[language=C, caption="Original grayscale code"]
	for (int y = 0; y < height; y++)
	{
		for (int x = 0; x < width; x++)
		{
			Vec3b color = color_image.at<Vec3b>(y, x);
			int gray = (color[0] + color[1] + color[2]) / 3;
			gray_image.at<unsigned char>(y, x) = gray;
		}
	}
\end{lstlisting}


\begin{lstlisting}[language=C, caption="CUDA kernel function to convert image to grayscale"]
	__global__ void convertToGray(uchar3 *input, unsigned char *output, int width, int height)
	{
		int x = blockIdx.x * blockDim.x + threadIdx.x;
		int y = blockIdx.y * blockDim.y + threadIdx.y;
	
		if (x < width && y < height)
		{
			output[y * width + x] = (input[y * width + x].x + input[y * width + x].y + input[y * width + x].z) / 3;
		}
	}
	
	dim3 block(BLOCK_SIZE, BLOCK_SIZE);
	dim3 grid((width + block.x - 1) / block.x, (height + block.y - 1) / block.y);
	
	convertToGray<<<grid, block>>>(d_color_image, d_gray_image, width, height);
\end{lstlisting}

\assignmentSection{Stage 4 – Analysis of results}

\paragraph*{MLP Analysis} As we can see, the OpenMP implementation is the fastest then the baseline and finally the CUDA and OpenMP is the slowest.
The CUDA profiler results for the mlp experiment are showed on table \ref{tab:mlp_cuda_profiler}.
The CUDA profiler shows that the CUDA memcpy is the most time consuming operation, followed by the backprop\_kernel and the CUDA memset.
So our kernel function is too short to benefit from the GPU acceleration. The cost to move the data to memory are too high. 
To conclude, we tried to optimize the application with OpenMP and CUDA. 
With OpenMP, we were able to speed up the application but not to great extent because of the hardware available. The CPU is only a quad core.
With CUDA, we wern't able to speed up the application because of the memory bandwidth bottlenecks described in the last section.

\paragraph*{Grayscale Analysis}
For the grayscale algorithm, we can see that the acceleration is substantial. For instance, for the winter image the GPU algorithm is 4 time faster. 
This is achieved with the grayscale kernel. The task is highly parallelizable and the kernel is able to use most of the threads available on the GPU
which is why the speedup is significant. 
As we can see in the results, the greater the image, the greater the speedup.

\begin{table}[h]\centering
	\begin{tabular}{|l|c|c|c|c|c|c|c|}
		\hline
		\textbf{Type} & \textbf{Time perc.} & \textbf{Time} & \textbf{Calls} & \textbf{Avg} & \textbf{Name}         \\ \hline
		GPU :         & 47.10\%             & 15.4163s      & 410754         & 37.531us     & CUDA memcpy HtoH      \\ \hline
		              & 35.72\%             & 11.6917s      & 68459          & 170.78us     & backprop\_kernel      \\ \hline
		              & 17.18\%             & 5.62353s      & 136918         & 41.072us     & CUDA memset           \\ \hline
		API calls:    & 74.08\%             & 145.662s      & 410754         & 354.62us     & cudaMemcpy            \\ \hline
		              & 13.21\%             & 25.9788s      & 136918         & 189.74us     & cudaMemset            \\ \hline
		              & 9.01\%              & 17.7185s      & 68459          & 258.82us     & cudaDeviceSynchronize \\ \hline
		              & 3.49\%              & 6.86282s      & 68459          & 100.25us     & cudaLaunchKernel      \\ \hline
		              & 0.18\%              & 344.80ms      & 7              & 49.257ms     & cudaMallocManaged     \\ \hline
		              & 0.04\%              & 72.091ms      & 68459          & 1.0530us     & cudaGetLastError      \\ \hline
		              & 0.00\%              & 953.77us      & 7              & 136.25us     & cudaFree              \\ \hline
		              & 0.00\%              & 108.55us      & 97             & 1.1190us     & cuDeviceGetAttribute  \\ \hline
		              & 0.00\%              & 10.521us      & 1              & 10.521us     & cuDeviceTotalMem      \\ \hline
		              & 0.00\%              & 7.3430us      & 3              & 2.4470us     & cuDeviceGetCount      \\ \hline
		              & 0.00\%              & 3.8540us      & 2              & 1.9270us     & cuDeviceGet           \\ \hline
		              & 0.00\%              & 1.5100us      & 1              & 1.5100us     & cuDeviceGetName       \\ \hline
		              & 0.00\%              & 938ns         & 1              & 938ns        & cuDeviceGetUuid       \\ \hline
	\end{tabular}
	\caption{MLP CUDA accelerated GPU profiling results}
	\label{tab:mlp_cuda_profiler}
\end{table}

\paragraph*{MLP future improvement}

Concerning MLP, with advanced knowledge it is surely possible to do the whole processing on GPU, which would surely be more optimized in the case of a nano jetson.
However we lack the knowledge to do it properly. 
Maybe there are other CUDA libraries that allow to do more complex operations and therefore to speed up some parts of the application that take a lot of time. 
I also think that the neural network is too small to be optimized properly and efficiently on GPU. 

\paragraph*{Grayscale future improvement}

Concerning the grayscale transforming application, there are already made functions of openCV optimized for the GPU that are even more efficient. 
However, this was not the purpose of this lab. 
There are even functions in OpenCV that integrate CUDA functions that allow to use GPU resources to the maximum.

\end{document}
